\documentclass[11pt]{article}
\usepackage[margin=1in]{geometry}
\usepackage{amsmath,amssymb,amsthm}
\usepackage{hyperref}
\newtheorem{definition}{Definition}
\newtheorem{lemma}{Lemma}
\newtheorem{proposition}{Proposition}
\newtheorem{corollary}{Corollary}
\theoremstyle{remark}
\newtheorem{remark}{Remark}
\newcommand{\E}{\mathbb{E}}
\newcommand{\I}{\mathcal{I}}
\newcommand{\Sset}{\mathcal{S}}
\newcommand{\Aset}{\mathcal{A}}

\title{RL for Sequential Recommendation:\\ Problem Formulation and First Results}
\author{Course project, topic 3.9}
\date{}

\begin{document}
\maketitle

\begin{abstract}
We cast supervised (short-term) and reinforcement-learning (long-term) recommenders as two solutions of a single
episodic combinatorial MDP with multinomial-logit (MNL) user choice, following Lee \& Oh (ICML 2025).
In this model the supervised next-click recommender is exactly the myopic ($\gamma = 0$) policy, the RL recommender
is the $\gamma > 0$ policy, and the slate size $K$ is an explicit parameter ($K = 1$ is classic sequential
recommendation, $K > 1$ is re-ranking). We prove two results that bound when RL can help:
a myopic policy can be arbitrarily worse than optimal on long-term return (a dual of Theorem 4.3 of
Labarca et al., ICML 2024), and RL has no advantage at all when the user's state transition does not depend on what
the user chose.
\end{abstract}

\section{Combinatorial MDP with MNL preference feedback}

\begin{definition}[Combinatorial MDP with MNL choice]
An episodic combinatorial MDP is a tuple $\mathcal{M} = (\Sset, \I, \Aset_K, \{P_h\}, \{r_h\}, \{p_h\}, H)$:
\begin{itemize}
  \item $\Sset$ is the set of user states (static features, recent history, and latent variables such as
  satisfaction or fatigue that the recommender may not observe);
  \item $\I = \{a_0, a_1, \dots, a_N\}$ is the item set, where $a_0$ is the outside option (``no click'');
  \item $\Aset_K = \{A \subseteq \I : a_0 \in A,\ 1 \le |A \setminus \{a_0\}| \le K\}$ is the set of slates;
  \item $P_h(a \mid s, A)$ is the probability that the user picks item $a \in A$;
  \item $p_h(s' \mid s, a)$ and $r_h(s, a)$ are the transition and reward after the user picks $a$;
  \item $H$ is the horizon (session length cap).
\end{itemize}
At step $h$ the recommender observes $s_h$, offers $A_h$, the user picks $a_h \sim P_h(\cdot \mid s_h, A_h)$,
the recommender receives $r_h(s_h, a_h)$ and the state moves to $s_{h+1} \sim p_h(\cdot \mid s_h, a_h)$.
\end{definition}

\paragraph{MNL choice (Lee \& Oh, Assumption 3.1).} With utilities $v(s,a) = \phi(s,a)^\top \theta$ and
$v(s, a_0) = 0$,
\begin{equation}
P(a \mid s, A) = \frac{\exp(v(s,a))}{1 + \sum_{a' \in A \setminus \{a_0\}} \exp(v(s,a'))}, \qquad a \in A.
\end{equation}
In our simulator the utility also depends on latent fatigue $f$:
$v(s,a) = \phi(s,a)^\top\theta - \lambda_f f \, \mathrm{sim}(a, \text{recent history})$.

\paragraph{Values.} For a policy $\pi = (\pi_1, \dots, \pi_H)$ and discount $\gamma \in [0,1]$,
the item-level and slate-level action values satisfy (Lee \& Oh, Sec.~3.1)
\begin{align}
\bar Q^\pi_h(s, a) &= r_h(s,a) + \gamma \sum_{s'} p_h(s' \mid s, a)\, V^\pi_{h+1}(s'), \\
Q^\pi_h(s, A) &= \sum_{a \in A} P_h(a \mid s, A)\, \bar Q^\pi_h(s, a), \qquad V^\pi_{H+1} \equiv 0.
\end{align}
The optimal policy maximises $V_1(s_1)$; the decomposition means only item-level values need to be learned,
and the slate is chosen by maximising a ratio of linear functions, which is a linear program
(Lee \& Oh, Appendix C) or, equivalently, a one-dimensional bisection.

\paragraph{Short-term and long-term metrics.} With $r = 1$ for a click and $0$ for $a_0$:
the short-term metric is the click-through rate of the first slate, $\mathrm{CTR}(\pi) = 1 - P(a_0 \mid s_1, \pi_1(s_1))$;
the long-term metric is the expected number of clicks in a session, $V^\pi_1(s_1)$ with $\gamma = 1$, together with
session length (steps before the user leaves).

\section{Supervised recommendation is the myopic policy}

\begin{definition}[Myopic policy]
The myopic policy chooses, in every state, $\pi^{\mathrm{my}}_h(s) \in \arg\max_{A \in \Aset_K}
\sum_{a \in A} P_h(a \mid s, A)\, r_h(s, a)$, i.e.\ the $\gamma = 0$ solution.
\end{definition}

\begin{lemma}[Supervised top-$K$ = myopic]\label{lem:sup}
Let rewards be $r(s,a) = 1$ for every item $a \neq a_0$ and $r(s, a_0) = 0$. Then the myopic slate is the set of
$K$ items with the largest utilities $v(s,a)$. In particular, a next-click model trained with cross-entropy that
recovers the MNL utilities (up to an additive constant) and recommends its top-$K$ scores is the myopic policy.
\end{lemma}
\begin{proof}
For a slate $A$ let $W(A) = \sum_{a \in A \setminus\{a_0\}} e^{v(s,a)}$. The expected immediate reward is
$\sum_{a \neq a_0} P(a \mid s, A) = W(A) / (1 + W(A))$, which is increasing in $W(A)$. Since every weight is
positive, $W$ is maximised over slates of at most $K$ items by the $K$ largest weights, i.e.\ the $K$ largest
utilities. The softmax next-click model $P(a \mid s) \propto e^{v(s,a)}$ is invariant to adding a constant to all
utilities, so its ranking equals the utility ranking.
\end{proof}

\begin{remark}
Lemma~\ref{lem:sup} is what lets us compare ``supervised vs RL'' inside one model: the topic's question becomes
``when is the $\gamma > 0$ optimal policy much better than the $\gamma = 0$ policy on long-term return?''
\end{remark}

\section{When RL can and cannot help}

\begin{proposition}[Myopic can be arbitrarily worse on long-term return]\label{prop:dual}
For every $n > 1$ there is a combinatorial MDP with MNL choice, unit click rewards and $\gamma = 1$ such that
$V^*_1(s_1) > n\, V^{\mathrm{my}}_1(s_1)$.
\end{proposition}
\begin{proof}
Take $\Sset = \{\mathrm{on}, \mathrm{off}\}$ with $\mathrm{off}$ absorbing and reward-free, $K = 1$, and two items:
a good item $g$ with utility $v_g$ and a junk item $j$ with utility $v_j > v_g$. Write
$c_x = e^{v_x}/(1 + e^{v_x})$ for the click probability of a slate holding item $x$, so $c_j > c_g$.
Clicking $j$ moves the user to $\mathrm{off}$ (the user leaves); clicking $g$ or not clicking keeps the user $\mathrm{on}$.
Rewards are 1 per click. By Lemma~\ref{lem:sup} the myopic policy offers $j$ in every step. Its return is the
number of clicks until the first click, which is at most 1, so $V^{\mathrm{my}}_1(\mathrm{on}) \le 1$
(precisely $1 - (1 - c_j)^H$). The policy that always offers $g$ never leaves $\mathrm{on}$ and collects
$H c_g$ in expectation, so $V^*_1(\mathrm{on}) \ge H c_g$. Choosing $H > n / c_g$ gives
$V^*_1 / V^{\mathrm{my}}_1 \ge H c_g > n$.
\end{proof}

\begin{remark}
Labarca et al.\ (Theorem 4.3) show the converse direction for one-step metrics: an optimal RL policy can be
arbitrarily worse than the cross-entropy model on next-item-prediction metrics. Together the two results say that
neither objective dominates; which recommender ``wins'' is decided by the evaluation metric and by the
dynamics, which motivates evaluating both short-term CTR and long-term return.
\end{remark}

\begin{proposition}[No advantage without choice-dependent dynamics]\label{prop:noadv}
Suppose that for every $h$ and $s$ the transition does not depend on the item chosen,
$p_h(\cdot \mid s, a) = p_h(\cdot \mid s)$ for all $a \in \I$ (including $a_0$). Then for every $\gamma \in [0,1]$
the myopic policy is optimal: $V^{\mathrm{my}}_1 = V^*_1$.
\end{proposition}
\begin{proof}
Let $c_h(s) = \gamma \sum_{s'} p_h(s' \mid s) V^*_{h+1}(s')$, which does not depend on the chosen item. Then for every
slate $A$,
\[
Q^*_h(s, A) = \sum_{a \in A} P_h(a \mid s, A) \bigl(r_h(s,a) + c_h(s)\bigr)
= \sum_{a \in A} P_h(a \mid s, A)\, r_h(s,a) + c_h(s),
\]
because the choice probabilities over $A$ (including $a_0$) sum to one. The second term is the same for every
slate, so $\arg\max_A Q^*_h(s, A)$ is exactly the myopic slate at every $(h, s)$. A policy that is greedy with
respect to $Q^*$ at every step is optimal, hence so is the myopic policy.
\end{proof}

\begin{corollary}
An RL recommender can improve long-term utility over a supervised one only if what the user picks changes the
distribution of the user's future state (satisfaction, fatigue, leaving). Measuring how strongly the
dynamics depend on the choice is therefore the first axis of our regime map.
\end{corollary}

\section{Slate size: a two-type toy model (to be analysed for the end-semester)}

Consider good items (utility $v_g$, the user stays after a click) and junk items (utility $v_j > v_g$,
the user leaves with probability $p$ after clicking one); with no click the user leaves with probability $p_0$.
A stationary slate holds $k$ junk and $K - k$ good items. Its click probability and per-step exit probability are
\[
c(k) = \frac{k e^{v_j} + (K-k) e^{v_g}}{1 + k e^{v_j} + (K-k) e^{v_g}}, \qquad
e(k) = p\, \frac{k e^{v_j}}{1 + k e^{v_j} + (K-k) e^{v_g}} + p_0\,(1 - c(k)),
\]
and its discounted long-term value is $V(k) = c(k) / \bigl(1 - \gamma(1 - e(k))\bigr)$.
The myopic policy uses the largest feasible $k$ (since $c(k)$ increases in $k$); the optimal stationary policy
maximises $V(k)$. The ratio $V(k^*)/V(k^{\mathrm{my}})$ as a function of $K$, $p$, $v_j - v_g$ and $\gamma$
is our prediction for how slate size changes the RL advantage; it will be compared with the simulator in the
second half of the project.

\section*{References}
\begin{itemize}
\item \'A.~Labarca, D.~Parra, R.~Toro Icarte. On the Unexpected Effectiveness of Reinforcement Learning for
Sequential Recommendation. ICML 2024.
\item J.~Lee, M.~Oh. Combinatorial Reinforcement Learning with Preference Feedback. ICML 2025.
\item Z.~Yuan, Q.~Wan, T.~Zhang, C.~Huo. PISDR: Page and Item Sequential Decision for Re-ranking Based on Offline
Reinforcement Learning. ACML 2024.
\item X.~Xin, A.~Karatzoglou, I.~Arapakis, J.~Jose. Self-Supervised Reinforcement Learning for Recommender
Systems (SQN). SIGIR 2020.
\end{itemize}
\end{document}
